% (C) 2026 Ralf Stephan, in collaboration with Claude Code.  CC0 1.0.
\chapter{Affine carry graphs}\label{ch:affine}

This chapter makes the recurrence~\eqref{eq:carry} a system in the sense of
Chapter~\ref{ch:covers} and builds its covers. Neither the system nor the graphs mention
certificates. A certificate of \texttt{DubC}, \texttt{DubC2}, \texttt{Oda26} or \texttt{Z32} is
a decidable witness for a property of one of these graphs. Matching each format to its graph
is the work of the bridges (Chapter~\ref{ch:splitting} for \texttt{DubC} and \texttt{Oda26},
the next milestone for \texttt{DubC2} and \texttt{Z32}).

\section{The affine system}

Fix integers \(p>q\ge1\) and a real shift \(\nu\). The \emph{affine step} is
\[
T(z)=\frac pq\,(z-\nu)+\nu,
\]
so \(z_n=\xi(p/q)^n+\nu\) satisfies \(z_{n+1}=T(z_n)\). The \emph{carry} of \(z\) is
\(s(z)=q\lfloor T(z)\rfloor-p\lfloor z\rfloor\).

\begin{lemma}\label{lem:affine}
For every real \(z\), with \(y=\frc z\) and \(y'=\frc{T(z)}\):
\begin{enumerate}[label=(\alph*)]
\item \(q\lfloor T(z)\rfloor=p\lfloor z\rfloor+s(z)\), the carry recurrence;
\item \(q\,y'=p\,y-(p-q)\nu-s(z)\), the fractional recurrence;
\item \(\lfloor-(p-q)\nu\rfloor-q\le s(z)<\lfloor-(p-q)\nu\rfloor-q+p+q+1\), the carry window.
\end{enumerate}
\end{lemma}

\begin{proof}
(a) is the definition. (b) follows from \(q\,T(z)=pz-(p-q)\nu\) and (a). For (c), write
\(s=p\,y-(p-q)\nu-q\,y'\) with \(y,y'\in[0,1)\) and compare with
\(\lfloor-(p-q)\nu\rfloor\le-(p-q)\nu<\lfloor-(p-q)\nu\rfloor+1\).
\end{proof}

Fix a window start \(s_0\) and a letter count \(n_A\). Letter \(t\in\{0,\dots,n_A-1\}\) stands
for the carry \(s_0+t\). For a finite set \(\mathcal P\) of positive integers (the divisor set),
\(x\) \emph{avoids} \(\mathcal P\) if no \(\ell\in\mathcal P\) divides \(x\).

\begin{definition}\label{def:affine-system}
The \emph{affine system}\index{affine system} is \((\R,T,\mathrm{adm})\) over the letters \(\{0,\dots,n_A-1\}\),
where \(\mathrm{adm}(z,t)\) holds iff \(s(z)=s_0+t\) and both \(\lfloor z\rfloor\) and
\(\lfloor T(z)\rfloor\) avoid \(\mathcal P\).
\end{definition}

Suppose the window contains all carries; Lemma~\ref{lem:affine}\,(c) gives one such window.
Then the integers \(\lfloor\xi(p/q)^n+\nu\rfloor\) avoid \(\mathcal P\) for all \(n\ge N\)
exactly when the orbit of \(z_N\) is admissible forever, with the carry word as its label
sequence.

\section{Residues, blocks and cells}

\begin{definition}[residue graph]\label{def:residue}
Let \(M\ge1\) and split the divisor set into \(P_M\) (divisors of \(M\)) and \(P_q\). The
\emph{residue graph}\index{residue graph} on \(\Z/M\) has an arc \(r\xrightarrow{t}r'\) iff
\(q\,r'=p\,r+(s_0+t)\) in \(\Z/M\), the residues \(r\) and \(r'\) avoid \(P_M\), and the carry
\(s_0+t\) avoids \(P_q\).
\end{definition}

\begin{proposition}\label{prop:cover-residue}
Suppose every \(\ell\in P_M\) divides \(M\) and every \(\ell\in P_q\) divides \(q\) and is
coprime to \(p\), with \(P_M\cup P_q\subseteq\mathcal P\). Then the residue graph covers the
affine system through \(\rho(r)=\{z:\lfloor z\rfloor\equiv r\pmod M\}\).
\end{proposition}

\begin{proof}
Reduce Lemma~\ref{lem:affine}\,(a) modulo \(M\). Divisibility by \(\ell\mid M\) depends only on
the residue. If some \(\ell\in P_q\) divided the carry
\(s=q\lfloor T(z)\rfloor-p\lfloor z\rfloor\), it would divide \(p\lfloor z\rfloor\), hence (being coprime to \(p\))
\(\lfloor z\rfloor\), which is avoided.
\end{proof}

The cell coordinate generalises to arbitrary intervals. Let \(I_i=[a_i,b_i)\)
(\(i=1,\dots,m\)) be intervals with rational ends. Lemma~\ref{lem:affine}\,(b) maps \(I_i\)
under the letter \(s\) to \(\bigl(p\,I_i-(p-q)\nu-s\bigr)/q\).

\begin{definition}[block and cell graphs]\label{def:block}
The \emph{block graph}\index{block graph} of \((I_i)\) at rational \(\nu\) has an arc \(i\xrightarrow{t}j\) iff the
image of \(I_i\) under the letter \(s=s_0+t\) meets \(I_j\):
\[
q\,a_j<p\,b_i-(p-q)\nu-s\quad\text{and}\quad p\,a_i-(p-q)\nu-s<q\,b_j .
\]
The \emph{cell graph}\index{cell graph} is the block graph of the \(D\) cells \([c/D,(c+1)/D)\)
(Example~\ref{ex:cell}). The \emph{affine graph}\index{affine graph} is the label product
\[
\mathrm{Aff}_{p,q,\nu}(M,D,P_M,P_q)=\mathrm{Res}_{p,q}(M,P_M,P_q)\otimes\mathrm{Cell}_{p,q,\nu}(D).
\]
\end{definition}

\begin{theorem}[the affine cover]\label{thm:cover-affine}\index{affine cover}
\begin{enumerate}[label=(\alph*)]
\item A block graph whose intervals cover \([0,1)\) covers the affine system through
  \(\rho(i)=\{z:\frc z\in I_i\}\). In particular the cell graph does.
\item Under the hypotheses of Proposition~\ref{prop:cover-residue}, the affine graph covers the
  affine system through \(\rho(r,c)=\{z:\lfloor z\rfloor\equiv r,\ \frc z\in[c/D,(c+1)/D)\}\),
  and every real \(z\) lies in some \(\rho(r,c)\).
\end{enumerate}
\end{theorem}

\begin{proof}
(a) The \emph{fractional system} has the maps \(y\mapsto(py-(p-q)\nu-s)/q\), and a step is
admissible when it stays in \([0,1)\). If \(y\in I_i\) and its image \(y'\) lies in \(I_j\),
then \(q\,a_j\le qy'<p\,b_i-(p-q)\nu-s\), and symmetrically, so the arc is present. By
Lemma~\ref{lem:affine}\,(b), \(z\mapsto\frc z\) maps admissible affine steps to admissible
fractional steps, and the cover pulls back (Chapter~\ref{ch:covers}, constructions).
(b) Proposition~\ref{prop:cover-residue}, part~(a) and the product principle.
\end{proof}

\paragraph{Gate G-4.} For the 16 cells of \texttt{DubC2/Atlas}, the successor lists of
\texttt{DubC2} (the function \lean{Params.succList}) were compared arc by arc, in exact rational arithmetic, with the
label product of Definition~\ref{def:block}: 30\,310 arcs, no mismatch. The \texttt{DubC2}
engine computes exactly \(\mathrm{Aff}\), including the rounding at cell boundaries.

\section{Realisability and refinement}

\begin{theorem}[Markov block graphs are exact]\label{thm:markov}\index{Markov block graph}
Let \(a_i<b_i\) and \(I_i\subseteq[0,1)\). Suppose every arc \(i\xrightarrow{t}j\) is
\emph{Markov}: \(I_j\) lies inside the image of \(I_i\). Then the block graph realises its
words in the fractional system. Every word from \(i\) is followed by all points of a nonempty
subinterval of \(I_i\).
\end{theorem}

\begin{proof}
Induction on the word, from the end. If the points of \([\alpha',\beta')\subseteq I_j\) follow
\(w\), the preimage of \([\alpha',\beta')\) under the letter \(t\) is a nonempty half-open
interval, and the Markov condition puts it inside \(I_i\).
\end{proof}

With Theorem~\ref{thm:count}\,(b), a right-resolving Markov block graph computes the entropy of
the carry language of the fractional system. The \texttt{Z32} type certificate is an instance.

Infinite paths are a different matter: nested half-open intervals can have an empty
intersection.

\begin{example}\label{ex:doubling}
The map \(y\mapsto 2y-s\) (\(p=2\), \(q=1\), \(\nu=0\)) on the single block \([0,1)\) with
letters \(s\in\{0,1\}\) is Markov, with equality at the right endpoint for \(s=1\). The constant
path of the letter \(1\) is not the itinerary of any point: \(1-y_n=2^n(1-y_0)\) leaves \([0,1]\).
Its only candidate is the fixed point \(1\notin[0,1)\).
\end{example}

\begin{theorem}[realisability forever]\label{thm:markov-inf}
If every arc \(i\xrightarrow{t}j\) is Markov with the right endpoint of \(I_j\) strictly inside the
image of \(I_i\), every infinite path of the block graph is the itinerary of a point of its first
interval that is admissible forever.
\end{theorem}

\begin{proof}
The points whose first \(m+1\) positions lie in the \emph{closed} intervals of the path form a
decreasing sequence of non-empty compact sets (non-empty by pulling a point back along the
path), so they have a common point \(y\). If \(y_{n+1}\le b_{j}\), the strict condition gives
\(y_n<b_i\), so every position lies in its half-open interval.
\end{proof}

\begin{proposition}[refinement]\label{prop:refine}
If each interval \(I_i\) lies inside \(J_{\pi(i)}\), then \(\pi\) is a homomorphism from the
block graph of \(I\) to the block graph of \(J\). In particular \(c\mapsto\lfloor c/k\rfloor\)
maps the cell graph with \(kD\) cells onto the one with \(D\) cells. By
Proposition~\ref{prop:labelproduct} and \lean{allPer\_of\_hom}, every label-form certificate of
the coarse graph holds for the fine one.
\end{proposition}

\section{The mirror}\label{sec:mirror}

Reflecting \([0,1)\) and negating the carries maps the affine system at \(\nu\) to the one at
\(1-\nu\). On the graph side this is exact.

\begin{proposition}\label{prop:mirror}
The arc condition of Definition~\ref{def:block} is invariant under
\((I,J,\nu,s)\mapsto(1-I,1-J,1-\nu,-s)\). Hence the cell graph at \(1-\nu\), with cells and letter
window reversed, has exactly the arcs of the cell graph at \(\nu\). If every prime of \(P_M\)
divides \(M\), the residue graph is invariant under \(r\mapsto-r\), \(s\mapsto-s\), and the affine
graphs at \(\nu\) and \(1-\nu\) are isomorphic up to relabelling. In particular one has the label
form iff the other has.
\end{proposition}

\begin{proof}
Write the arc condition out: the two inequalities exchange roles. Divisibility of \(r\) and of
\(M-r\) by \(\ell\mid M\) agree.
\end{proof}

The half-open cells reflect to cells open on the other side, so the covers are not mirror
images at the endpoints. Each graph covers its own system (Theorem~\ref{thm:cover-affine}),
and that is all Theorem~\ref{thm:divisor-any-cover} below needs: a certified cell at \(\nu\) is
a theorem at \(1-\nu\), with no orbit-level argument.

\section{Parameter intervals}

The arc condition of Definition~\ref{def:block} is linear in \(\nu\). With
\(\alpha=p\,b_i-s-q\,a_j\) and \(\beta=p\,a_i-s-q\,b_j\) it reads \(\beta<(p-q)\nu<\alpha\).

\begin{definition}\label{def:hull}
For rational \(\nu_0\le\nu_1\), the \emph{hull}\index{hull} block graph has an arc \(i\xrightarrow{t}j\) iff
\((p-q)\nu_0<\alpha\), \(\beta<(p-q)\nu_1\) and \(\beta<\alpha\). The \emph{affine hull}\index{hull!affine} is the
product of the residue graph with the hull of the cells.
\end{definition}

\begin{theorem}\label{thm:hull}
Let \(q<p\).
\begin{enumerate}[label=(\alph*)]
\item The hull is the union of the block graphs at rational \(\nu\in[\nu_0,\nu_1]\).
\item For every real \(\nu\in[\nu_0,\nu_1]\), the hull covers the affine system at \(\nu\), and
  so does the affine hull.
\item A smaller interval gives a subgraph of the hull. A larger divisor set gives a subgraph of
  the residue graph, and the label form of the affine graph passes from \((P_M,P_q)\) to every
  \((P'_M,P'_q)\supseteq(P_M,P_q)\).
\item The window \(s_0\le\lfloor-(p-q)\nu_1\rfloor-q\),
  \(\lfloor-(p-q)\nu_0\rfloor+p+1\le s_0+n_A\) contains every carry for every
  \(\nu\in[\nu_0,\nu_1]\).
\end{enumerate}
\end{theorem}

\begin{proof}
(a) Take \((p-q)\nu\) halfway between \(\max((p-q)\nu_0,\beta)\) and
\(\min((p-q)\nu_1,\alpha)\). (b) The arc condition at real \(\nu\) implies the three hull
inequalities. Then apply Theorem~\ref{thm:cover-affine}. (c) Homomorphisms by the identity.
(d) Lemma~\ref{lem:affine}\,(c) and monotonicity of the floor.
\end{proof}

Part~(b) is the mechanism of \texttt{DubC2/NuAtlas} and \lean{DubC2.Params.hullOK}: one
certificate on a hull settles a whole interval of real shifts at once.

\section{The dimension of the avoiders}

\begin{theorem}\label{thm:dim-avoiders}
Let \(q<p\), and let \(X\) be any finite carry graph that covers the affine system through
\(\rho\), with \(\bigcup_v\rho(v)=\R\) and a window containing every carry. Then
\[
\dim_H\bigl\{\xi\in\R:\ \lfloor\xi(p/q)^n+\nu\rfloor\text{ avoids }\mathcal P\text{ for all
large }n\bigr\}\ \le\ \max\Bigl(0,\frac{h(X)}{\log(p/q)}\Bigr).
\]
In particular this holds for the affine graph and, uniformly in \(\nu\in[\nu_0,\nu_1]\), for the
affine hull.
\end{theorem}

\begin{proof}
Split the set by the first time \(N\) from which it avoids \(\mathcal P\) and by the integer
part \(x_0\) of \(z_N\). The point \(z_N\) is admissible forever. Two such points with the same
integer part that follow the same word keep equal integer parts along the run, by the carry
recurrence. Their runs therefore stay within distance \(1\), and \(T\) expands by exactly
\(p/q\). So the itinerary localises with ratio \(q/p\), and Theorem~\ref{thm:dimension} bounds
each piece. The map \(z_N\mapsto\xi\) is Lipschitz, and a countable union does not raise the
Hausdorff dimension.
\end{proof}

This is \lean{DubC2.Gen.dimH\_eventually\_avoids\_le} with the certificate removed: the graph
can be any cover, and the bound is its entropy. A super-eigenvector bracket
(Theorem~\ref{thm:brackets}) turns it into the explicit number \(\log(a/b)/\log(p/q)\). At
\(\mathcal P=\{2\}\) and \(p/q=3/2\) the bound is \(1\): there the model has entropy
\(\log(3/2)\), which is Mahler's problem \cite{Mah68}.

\section{The certificates as carry graphs}\label{sec:cert-bridges}

The two certificate roots of the corpus check finite witnesses for the graphs of this chapter.
The bridge files of the root \texttt{CarryBridge/} spell out the dictionary. They re-derive the roots' soundness theorems from the generic ones, using one
arithmetic fact:

\begin{theorem}[the divisor theorem through any cover]\label{thm:divisor-any-cover}\index{divisor theorem}
Let \(1<q<p\) be coprime, and let \(X\) be any carry graph that covers the affine system through
\(\rho\), with \(\bigcup_v\rho(v)=\R\) and a window containing every carry. If \(X\) has the
label form, then for every \(\xi\ne0\) infinitely many of \(\lfloor\xi(p/q)^n+\nu\rfloor\) are
divisible by a member of \(\mathcal P\).
\end{theorem}

\begin{proof}
Otherwise some tail \(z_N=\xi(p/q)^N+\nu\) is admissible forever. By
Theorem~\ref{thm:transfer-per} its carry sequence is eventually periodic. The carry word of
\(\xi(p/q)^n+\nu\) is not eventually periodic for \(\xi\ne0\) \cite[Lemma~2]{DN05}.
\end{proof}

\paragraph{DubC2.} A cell \(\mathcal C=(p,q,a/b,M,D,P_M,P_q,s_0,n_A)\) of \texttt{DubC2.Gen} lists
the arcs of the affine graph \(\mathrm{affine}(p,q,a/b,s_0,n_A,M,D,P_M,P_q)\), with the vertex
\((r,c)\) numbered \(rD+c\). Multiplying the arc condition of Definition~\ref{def:block} by
\(bD\) gives the certificate's integer test
\[
qb\,c'<pb(c+1)-K(s),\qquad pb\,c-K(s)<qb(c'+1),
\]
where \(K(s)=(p-q)aD+sbD\).
The candidate range \([c_{\mathrm{lo}},c_{\mathrm{hi}})\) of \lean{succList} contains every
\(c'\) that passes this test, and the residue step \(r'=(pr+s)q^{-1}\bmod M\) is the residue
graph's. So the certificate graph and the affine graph are isomorphic. A checked certificate
is a witness for the label form of the affine graph in exactly the shape of
Theorem~\ref{thm:peel-simulate}: its rounds are a peel chain (Definition~\ref{def:peel-chain})
ending inside a listed core, and its table of subset states is a carry graph with a labelled
rank that simulates the certificate graph on the core. Theorem~\ref{thm:divisor-any-cover} then
gives \texttt{DubC2}'s divisor theorem, and by Proposition~\ref{prop:mirror} the same theorem at
the mirror shift \(1-a/b\).

The other DubC2 witnesses are read the same way:
\begin{itemize}
\item A hull check over \([a_0/b,a_1/b]\) is a homomorphism from the affine hull into the graph
at \(a/b\). By Theorem~\ref{thm:hull}, one certificate then settles every real shift of the
interval.
\item An entropy certificate is a subset automaton that simulates the affine graph, so it covers
the system as well. Its weights form a super\-/eigenvector, and
Theorem~\ref{thm:dim-avoiders} gives the bound \(\log(a/b)/\log(p/q)\).
\item A cycle certificate lists the cycles of the rank-keeping part of the subset automaton
and a potential. With rank \(0\) on the listed states it is a peel rank, so every cycle of the
rank-keeping part is listed (Theorem~\ref{thm:peel}).
\item The follower automaton is a right-resolving carry graph on an infinite vertex set. ``Some
\(D\) certifies'' becomes: finitely many vertices are reachable, and none of them branches. One
direction is geodesic Kőnig plus one arithmetic lemma: a periodic letter sequence cannot drive
an injective follower path. A follower tree check is a labelled rank on the reachable part.
\item A horseshoe witness yields an avoiding pseudo-orbit whose letters are not eventually
periodic, so no cover has the label form. A branching witness only yields finite pseudo-orbit
segments, but every prefix of an aperiodic word is then read, and Corollary~\ref{cor:konig-labels}
gives an infinite path with aperiodic labels.
\end{itemize}

\paragraph{Z32.} A block certificate at base \(p/q\) (with \(\nu=0\)) gives a block graph whose
vertices are the distinct blocks and whose arcs are the \texttt{hits} relation. Its
stratification is a labelled rank. The funnel puts every orbit that stays in \(U\) on an
infinite path of this graph, so the rank theorem and \cite[Lemma~2]{DN05} rule such orbits out.
The other \texttt{Z32} results are statements about finite functional carry graphs as well:
\begin{itemize}
\item The escape bound needs a time after which the carry word repeats. For an unranked
certificate it is Proposition~\ref{prop:functional-walk} on the block graph, with budget
\(|V|\le\#\text{blocks}\).
\item The depth-\(K\) schema of the windows \([s,s+1/p)\) is the rotation of \(K+2\) ranks by the
low and the high branch, a functional graph on \(K+2\) vertices. A confined orbit's ranks walk
it, which gives both the qualitative theorem and the budget \(K+2\).
\item Two confined orbits in one block of an unranked certificate read the same carry word,
because a path in a functional graph is determined by its start. Equal words separate at the
rate \((p/q)^n\), so the hold set is finite and the removed set must meet all but finitely many
cycles.
\item A type certificate for an anchored window \([0,\ell)\) is a right-resolving carry graph that
covers and realises the window's model system exactly. Its two weight vectors are a super- and a
sub-eigenvector (Theorem~\ref{thm:brackets}), and its label language from the root is the
window's language.
\end{itemize}

\begin{example}\label{ex:two-cell}
Take \(p/q=3/2\) and \(U=[0,1/5)\cup[4/5,1)\). The certificate has denominator \(15\) and funnel
depth \(1\), and its two blocks are \(B_0=[0,2/15)\) and \(B_1=[13/15,1)\). Among the four carries
\(s\in\{-1,0,1,2\}\) the block graph has exactly two arcs, both loops: \(B_0\xrightarrow{0}B_0\)
and \(B_1\xrightarrow{1}B_1\). It is functional. A confined orbit would therefore have carry word
eventually \(0^\infty\) or \(1^\infty\), which is impossible for \(\xi\ne0\). Hence no \(\xi\ne0\)
has \(\lVert\xi(3/2)^n\rVert<1/5\) for all \(n\) (with the half-open convention at \(4/5\)).
\end{example}

A horseshoe certificate for the same \(U\) would realise \(N^k\ge2^k\) model words of length
\(kL\). A functional block graph allows at most one word per block, so the two kinds of
certificate never meet, by counting alone (\lean{not\_cert\_and\_horse'}).

\section*{Lean}

\begin{center}\small
\begin{tabular}{@{}>{\raggedright}p{0.3\textwidth}>{\raggedright\arraybackslash}p{0.64\textwidth}@{}}
\toprule
Text & Lean (\texttt{CarryGraph/Affine.lean}, \texttt{Interval.lean})\\
\midrule
Lemma~\ref{lem:affine} & \lean{affStep}, \lean{carry}, \lean{affStep\_orbit},
  \lean{q\_mul\_floor\_affStep}, \lean{q\_mul\_affStep}, \lean{fract\_affStep},
  \lean{carry\_window}\\
Definition~\ref{def:affine-system} & \lean{Avoid}, \lean{affAdm}\\
Definition~\ref{def:residue}, Proposition~\ref{prop:cover-residue} & \lean{ResOK},
  \lean{residueGraph}, \lean{adj\_residueGraph}, \lean{resOK\_intCast},
  \lean{cover\_residueGraph}\\
Definition~\ref{def:block} & \lean{BlockArc}, \lean{ival}, \lean{blockGraph},
  \lean{adj\_blockGraph}, \lean{cells}, \lean{cellGraph}, \lean{affine}\\
Theorem~\ref{thm:cover-affine} & (a) \lean{fracStep}, \lean{fracAdm},
  \lean{cover\_frac\_of\_blockArc}, \lean{fract\_semiconj}, \lean{cover\_blockGraph},
  \lean{exists\_cell}, \lean{cover\_cellGraph}; (b) \lean{cover\_affine},
  \lean{affine\_rho\_cover}\\
Theorem~\ref{thm:markov} & \lean{realizes\_blockGraph}\\
Example~\ref{ex:doubling}, Theorem~\ref{thm:markov-inf} & \lean{not\_realizesInf\_doubling},
  \lean{continuous\_orbit\_fracStep}, \lean{realizesInf\_blockGraph}\\
Proposition~\ref{prop:mirror} & \lean{blockArc\_mirror}, \lean{cells\_rev},
  \lean{adj\_cellGraph\_mirror}, \lean{resOK\_neg}, \lean{adj\_residueGraph\_mirror},
  \lean{adj\_affine\_mirror}, \lean{allPer\_affine\_mirror\_iff}, \lean{allPer\_iff\_of\_adj\_equiv}
  (\texttt{CarryGraph/Mirror.lean})\\
Proposition~\ref{prop:refine} & \lean{blockGraph\_hom}, \lean{cellGraph\_hom}\\
Definition~\ref{def:hull}, Theorem~\ref{thm:hull} & \lean{blockHull}, \lean{adj\_blockHull},
  \lean{cellHull}, \lean{affineHull}; (a) \lean{adj\_blockHull\_iff}; (b)
  \lean{blockArc\_hull}, \lean{cover\_blockHull}, \lean{cover\_affineHull}; (c)
  \lean{blockHull\_mono}, \lean{blockHull\_hom}, \lean{residueGraph\_hom},
  \lean{allPer\_affine\_mono}; (d) \lean{carry\_window\_of\_interval}\\
Theorem~\ref{thm:dim-avoiders} & \lean{floor\_run\_eq}, \lean{survivors},
  \lean{dimH\_survivors\_le}, \lean{dimH\_avoiders\_le}, \lean{dimH\_avoiders\_le\_hull}\\
Theorem~\ref{thm:divisor-any-cover} & \lean{CarryBridge.frequently\_not\_avoid\_of\_allPer}
  (\texttt{CarryBridge/DubC2.lean})\\
Section~\ref{sec:cert-bridges} & \texttt{CarryBridge/}: \texttt{DubC2}, \texttt{DubC2Witness},
  \texttt{Z32}, \texttt{Z32Cert}, \texttt{Z32Type}, \texttt{Z32Balance}; see
  Appendix~\ref{app:dictionary}\\
\bottomrule
\end{tabular}
\end{center}
